\par\Needspace{6\baselineskip}\medskip
\begin{center}\hypertarget{native-199-article149}{}\pdfbookmark[2]{149.}{article-149}149.\end{center}
Only in the last hypothesis are the elements themselves calculated, either from $f,r',r''$ or from $f'',r,r'$, by completing one of the two calculations that in the preceding hypotheses needed to be taken only as far as $\eta$ or $\eta''$. If we choose to complete both, the agreement of the resulting numbers will provide a further check on the entire work. It is preferable, however, as soon as $f,f',f''$ have been found, to derive the elements from the combination of the first and third positions alone, that is, from $f',r,r''$ and the time interval. Finally, as a further check on the calculation, determine the middle position in the orbit from the elements found.

In this way we obtain the dimensions of the conic section, namely the eccentricity, the semimajor axis or the semiparameter, and the position of perihelion relative to the heliocentric positions $C,C',C''$, as well as the mean motion and the mean anomaly for an arbitrary epoch if the orbit is elliptical, or the time of passage through perihelion if the orbit is hyperbolic or parabolic. It remains to determine the positions of the heliocentric points in the orbit relative to the ascending node, the position of that node relative to the equinoctial point, and the inclination of the orbit to the ecliptic (or equator). All this can be accomplished by solving a single spherical triangle. Let $\AscendingNode$ be the longitude of the ascending node, $i$ the inclination of the orbit, and $g$ and $g''$ the arguments of latitude in the first and third observations. Finally, let $l-\AscendingNode=h$, $l''-\AscendingNode=h''$. With $\AscendingNode$ denoting the ascending node in the fourth figure, the sides of triangle $\AscendingNode AC$ are $AD'-\zeta,g,h$, and the angles respectively opposite them are $i,180^{0}-\gamma,u$.
Thus we have
\SourceLeaf{200}
\[
\begin{aligned}
\sin\tfrac12i\sin\tfrac12(g+h)&=\sin\tfrac12(AD'-\zeta)\sin\tfrac12(\gamma+u)\\
\sin\tfrac12i\cos\tfrac12(g+h)&=\cos\tfrac12(AD'-\zeta)\sin\tfrac12(\gamma-u)\\
\cos\tfrac12i\sin\tfrac12(g-h)&=\sin\tfrac12(AD'-\zeta)\cos\tfrac12(\gamma+u)\\
\cos\tfrac12i\cos\tfrac12(g-h)&=\cos\tfrac12(AD'-\zeta)\cos\tfrac12(\gamma-u).
\end{aligned}
\]

The first two equations give $\tfrac12(g+h)$ and $\sin\tfrac12i$, and the remaining two $\tfrac12(g-h)$ and $\cos\tfrac12i$. From $g$ we obtain the position of perihelion relative to the ascending node, and from $h$ the position of the node in the ecliptic. Finally we obtain $i$, with the sine and cosine providing a check on each other. We can achieve the same result with the triangle $\AscendingNode A''C''$, for which we need only change the symbols $g,h,A,\zeta,\gamma,u$ to $g'',h'',A'',\zeta'',\gamma'',u''$ in the preceding formulas. As a further check on the entire work it is worth carrying out the calculation in both ways. If there are any slight differences between the values of $i,\AscendingNode$ and of the longitude of perihelion in the orbit, it is appropriate to take the mean values. These differences will seldom reach $0''{,}1$ or $0''{,}2$ if all the calculations have been carried out accurately to seven decimal places.

Whenever the equator is adopted as the fundamental plane in place of the ecliptic, the sole difference in the calculation is that the points $A,A''$ must be replaced by the intersections of the equator with the great circles $AB,A''B''$.
